% SOURCE: Chinese PDF viewer pp. 97--98. Section 1.
\section{The Idea of Stationary Perturbation Theory}
The Schr\"odinger equation is
\begin{equation}
 \ii\hbar\frac{\partial}{\partial t}\ket{\psi,t}
 =\op H\ket{\psi,t}.
\end{equation}
Only a very small number of physical problems can be solved exactly. For most actual problems, the construction of $\op H$ is so complicated that an exact solution is impossible; approximate methods must be used. Different problems call for different methods of approximation.

When $\op H$ does not depend on time, separation of variables gives a stationary state
\begin{equation}
 \ket{\psi_n,t}=\ket n e^{-\ii E_nt/\hbar},
\end{equation}
where the energy $E_n$ and the energy eigenstate $\ket n$ satisfy the stationary Schr\"odinger equation
\begin{equation}
 \op H\ket n=E_n\ket n.
\end{equation}
Write
\begin{equation}
 \op H=\op H^{(0)}+\op H^{(1)}.
\end{equation}
The following are required:
\begin{enumerate}
\item $\op H^{(0)}$ contains the main part of $\op H$, so $\op H^{(1)}$ is small. Since both $\op H^{(0)}$ and $\op H^{(1)}$ are operators, their relative magnitudes can be understood from the classical correspondence principle. Strictly speaking, one compares the matrix elements of the two operators, as discussed below.
\item The stationary equation for $\op H^{(0)}$,
\begin{equation}
 \op H^{(0)}\ket n^{(0)}=E_n^{(0)}\ket n^{(0)},
\end{equation}
can be solved exactly.
\end{enumerate}
Expand $E_n$ and $\ket n$:
\begin{equation}
 E_n=E_n^{(0)}+E_n^{(1)}+E_n^{(2)}+\cdots,
 \qquad
 \ket n=\ket n^{(0)}+\ket n^{(1)}+\ket n^{(2)}+\cdots.
\end{equation}
Insert these into the stationary Schr\"odinger equation:
\begin{equation}\begin{aligned}
 &(\op H^{(0)}+\op H^{(1)})
 (\ket n^{(0)}+\ket n^{(1)}+\ket n^{(2)}+\cdots)\\
 &\hspace{10mm}=(E_n^{(0)}+E_n^{(1)}+E_n^{(2)}+\cdots)
 (\ket n^{(0)}+\ket n^{(1)}+\ket n^{(2)}+\cdots).
\end{aligned}\end{equation}
The equation is then solved successively, order by order.

\textbf{Zeroth-order approximation:}
\begin{equation}
 \op H^{(0)}\ket n^{(0)}=E_n^{(0)}\ket n^{(0)}
 \quad\longrightarrow\quad E_n^{(0)},\ \ket n^{(0)}.
\end{equation}
\textbf{First-order approximation:}
\begin{equation}
 \op H^{(0)}\ket n^{(1)}+\op H^{(1)}\ket n^{(0)}
 =E_n^{(0)}\ket n^{(1)}+E_n^{(1)}\ket n^{(0)},
\end{equation}
or
\begin{equation}
 (\op H^{(0)}-E_n^{(0)})\ket n^{(1)}
 =-(\op H^{(1)}-E_n^{(1)})\ket n^{(0)}
 \quad\longrightarrow\quad E_n^{(1)},\ \ket n^{(1)}.
\end{equation}
\textbf{Second-order approximation:}
\begin{equation}
 (\op H^{(0)}-E_n^{(0)})\ket n^{(2)}
 =-(\op H^{(1)}-E_n^{(1)})\ket n^{(1)}
 +E_n^{(2)}\ket n^{(0)}
 \quad\longrightarrow\quad E_n^{(2)},\ \ket n^{(2)}.
\end{equation}
And so forth. In general, only the first nonzero correction is calculated.

Perturbation theory is based on the eigenvalues $E_n^{(0)}$ and eigenstates $\ket n^{(0)}$ of $\op H^{(0)}$. In obtaining higher-order corrections, it is necessary to consider whether $E_n^{(0)}$ is degenerate. We first consider the nondegenerate case.